\subsection*{Chapter 24: Fine-Tuning for Classification}

\paragraph{Exercise 24.1.} With global padding every batch is $8\times120$; with dynamic
padding the mean batch width falls to roughly the batch's longest message (typically
30--60 tokens for SMS), so epoch time drops roughly in proportion once the GPU is
compute-bound. On CPU, or with tiny batches, the gain is smaller because fixed overheads
dominate. Accuracy should be statistically indistinguishable \emph{if} the head reads
the last real token. Padding to 1{,}024 multiplies linear-layer work by about
$1024/120\approx8.5$ and attention-score work by about $73$. If the head reads position
$-1$, it now reads a padding token that sits 900+ positions after the message. Its
hidden state is still a function of the message, but a very different input
distribution from the one seen in training, and accuracy typically drops sharply.

\paragraph{Exercise 24.2.} Per-step cost is predicted by $2N + 4N_{\text{tr}}$ per token:
with $k$ trainable blocks out of 12, $N_{\text{tr}}\approx (k/12)N$, so time per step
$\propto 2 + 4k/12$. That is $1.0$ ($k=0$, head only, up to a tiny head), $1.17$
($k=1$), $1.33$ ($k=2$), $1.67$ ($k=4$) and $3.0$ ($k=12$) relative to a forward pass.
Peak memory grows with $k$ through activations kept for backward (only blocks above the
lowest trainable one store them) and through $14$ bytes per trainable parameter.
Accuracy usually improves from $k=0$ to $k=1$ and changes little after that.
Discrepancies come from fixed overheads (optimizer step, kernel launches) and from the
embedding/head layers, which the block-count model ignores.

\paragraph{Exercise 24.3.} Run blocks 1--11 under \code{torch.no\_grad()} once, storing
$h_{11}$ of shape $(n, T, 768)$ in bf16 (about \SI{193}{MB} for 1{,}045 messages padded
to 120). Train block 12, the final norm and the head on the cache. With dropout off,
the forward computation is identical, so the loss curves match to floating-point
tolerance. Measuring cost in units of one block's forward FLOPs (2 per parameter), an uncached step costs $12\times2 + 4 = 28$ (twelve forwards, one backward of block 12), a cached step $2+4=6$, so each epoch is about $28/6\approx4.7\times$ faster (less if data movement of the cache dominates).

\paragraph{Exercise 24.4.} Plain cross-entropy on 87/13 data reaches high accuracy but
lower spam recall. Class weights $w_c\propto 1/\text{freq}(c)$ raise recall at some cost
in precision. Undersampling behaves similarly but trains on less data. Accuracy is
misleading under imbalance; report precision/recall/F1 for the spam class, and for a
dashboard prefer precision at a fixed recall (or vice versa), matched to the product
cost of each error.

\paragraph{Exercise 24.5.} Let a message occupy positions $0..\ell-1$ with padding at
$\ell..T-1$. Under the causal mask, the hidden state at position $t$ depends only on
tokens $0..t$. In particular position $\ell-1$ never attends to padding, and
LayerNorm/MLP act per position, so $h_{\ell-1}$, and hence the logits read there, are
independent of $T-\ell$. Counterexample for position $-1$: with $T=\ell+1$, the last
position is a padding token attending to $\ell+1$ keys. With $T=\ell+5$ it attends to
$\ell+5$ keys, four of which are padding tokens with different position embeddings,
so the softmax weights and the output differ.

\paragraph{Exercise 24.6.} With left padding the real tokens occupy positions
$T-\ell..T-1$ instead of $0..\ell-1$, so they receive different learned position
embeddings, and the logits differ (often by a few percent of their magnitude, enough to
flip borderline cases). The key mask removes the influence of padded tokens but cannot
undo the position shift. With RoPE, attention depends only on $i-j$, which is unchanged
for every pair of real tokens, so the logits match up to floating-point error.

\paragraph{Exercise 24.7.} Sort-within-mega-batch as in \cref{lst:classify:bucket}, then
greedily add sequences to the current batch while $\max(\text{len})\times(\text{count})
\le K$. Short-message batches then contain many sequences and long ones few. The number
of batches varies, padding waste matches bucketing, and step-time variance drops
because every batch has roughly $K$ token-slots. Learning-rate schedules should then be
defined in tokens rather than steps.

\paragraph{Exercise 24.8.} Compute $\log p(\text{`` spam''}\mid\text{prompt})$ and
$\log p(\text{`` ham''}\mid\text{prompt})$ from one forward pass each (or one pass with
both continuations scored via \code{sequence\_logprob}) and predict the larger. Base
GPT-2 models typically score only modestly above chance on the balanced split, and the
improvement with model size is gradual. Fine-tuning closes most of the gap. Note that
the label words should be chosen as single tokens with a leading space, otherwise
multi-token labels are penalised by length.

\paragraph{Exercise 24.9.} Demand: $20{,}000\times35 = 7.0\times10^5$ real tokens/s. With
random batching at about 60\% waste, the GPUs must process $7.0\times10^5/0.4=1.75\times
10^6$ token-slots/s; with bucketing at about 7\% waste, $7.5\times10^5$. An L4 offers
\SI{121}{TFLOP/s}; at 40\% utilisation and $1.7\times10^8$ FLOPs/token it processes
about $2.85\times10^5$ tokens/s. So random batching needs $\lceil 6.1\rceil = 7$ GPUs
and bucketing needs $\lceil 2.6\rceil=3$, plus headroom for peaks. A 160-token message
in a batch of 64 takes about $64\times160\times1.7\times10^8/4.8\times10^{13}\approx
\SI{36}{ms}$ of compute, so the \SI{50}{ms} p99 constrains the batching window to a few
milliseconds. The weakest assumption is the 40\% utilisation: at such short sequences,
kernel-launch and framework overheads can halve it unless CUDA graphs or a compiled
engine are used.

\paragraph{Exercise 24.10.} The SMS collection contains several hundred exact duplicate
texts, mostly spam templates, and more after normalisation. Re-splitting after
deduplication typically lowers test accuracy slightly. The difference measures how much
the original score benefited from test messages whose duplicates were in training.

\paragraph{Exercise 24.11.} Temperature scaling multiplies all logits of an example by the
same positive constant $1/T$. Because $\argmax_c z_c/T = \argmax_c z_c$ for $T>0$,
predictions and therefore accuracy are unchanged. A fine-tuned classifier is typically
over-confident ($T>1$ after fitting), and ECE drops substantially after scaling.

\paragraph{Exercise 24.12.} Sort test examples by calibrated confidence. Escalating the
least confident fraction $e$ makes cascade accuracy
$=1-(\text{errors among the confident }1-e)/n$. Increase $e$ until accuracy reaches
99.5\%; the cost is then $(1+eR)/R$ of the large-only budget. Well-calibrated models
concentrate their errors in the low-confidence region, so a small $e$ suffices.

\paragraph{Exercise 24.13.} For pretrained GPT-2 small, probe accuracy typically rises over
the first several blocks and plateaus in the middle-to-upper layers. Sometimes it dips
slightly in the last block, which is specialised for next-token prediction. If the
plateau is reached by block $\ell$, serving only blocks $1..\ell$ plus a head saves
about $(12-\ell)/12$ of the per-token FLOPs (embeddings and head aside).
